\documentclass[aspectratio=1610]{beamer}
\mode<presentation>
\usetheme{Madrid}
\usecolortheme{seahorse}
% \usepackage[table]{xcolor}
\usepackage{multimedia}
\usepackage{media9}
\usepackage[utf8]{inputenc}
\usepackage{amsmath}
\usepackage{lmodern}
% \usepackage[usenames, dvipsnames]{color}

\usepackage{arydshln}
\usepackage{url}
\usepackage[absolute,overlay]{textpos}
\setlength{\TPHorizModule}{\textwidth}
\setlength{\TPVertModule}{\textwidth}

\graphicspath{{plots/}}

\setbeamercolor{CBwonb}{fg=white,bg=white!10!blue}
\setbeamercolor{CBronb}{fg=red,bg=white!10!blue}

\usepackage{tikz}
\usetikzlibrary{shapes,arrows,positioning}
\tikzstyle{startstop} = [rectangle, rounded corners, minimum width=1cm, minimum height=0.1cm,text centered, text width=2.5cm, draw=black, fill=red!30]
\tikzstyle{io} = [trapezium, trapezium left angle=70, trapezium right angle=110, minimum width=1cm, minimum height=0.1cm, text centered, text width=2.5cm, draw=black, fill=blue!30]
\tikzstyle{process} = [rectangle, minimum width=1cm, minimum height=0.1cm, text centered,text width=2.5cm, draw=black, fill=orange!30]
\tikzstyle{decision} = [circle, minimum width=1cm, minimum height=0.1cm, text centered,text width=2.5cm, draw=black, fill=green!30]
\tikzstyle{arrow} = [thick,->,text centered, color=blue, text width=2cm,>=stealth]

\newcommand{\Msun}{~{\rm M_\sun}}
\newcommand{\hMsun}{~h^{-1}\>{\rm M_\odot}}
\newcommand{\Mpc}{~h^{-1}~{\rm Mpc}}
\newcommand{\Kpc}{~h^{-1}~{\rm kpc}}


\title[]{The Large Scale Environments}
\subtitle{The large-scale distribution of baryons inside the cosmological
hydrodynamical simualtions.}
\author[Email: weiguang.cui@uam.es]{{\Large \bf Weiguang Cui},\inst{*} \footnote{\url{https://weiguangcui.github.io}}}
\institute[]{
  \inst{*}
  Departamento de F\'isica Te\'{o}rica, \\
  Universidad Aut\'{o}noma de Madrid, 28049 Madrid, Spain
}
\date[]{The HUBS workshop, Shanghai. \\  October, 16, 2018}
\logo{\includegraphics[height=1cm]{logo_uam.jpg}}

%%%%%%%%%%%%%%%%%%%%%%%%%%%%%%%%%%%%%%%%%%%%%%%%%%%%%%%%%%%%%%%%%%%%%%%%%
\begin{document}
  \frame{\titlepage}

%----------------------------------
\section{Introduction} \label{sec:1}

\begin{frame}
  \frametitle{Background: Cluster of Galaxies:}
    \begin{itemize}
      \item<1-> Galaxy cluster is the final state of the hierarchical structure formation.
      \item<2-> Typically a few Mpc across, contain $\sim$ 100 - 1000 luminous galaxies and many more dwarfs, masses $\gtrsim 10^{14.5} M_{\odot}$.
      \item<3-> Dark matter: $\sim 85 \%$, gas: $\sim 13 \%$, star: $\sim 1 - 2 \%$.
      \item<4-> It provides plentiful information for baryonic processes, galaxy formation and cosmology models.
    %   \item<5->[]
    %     \begin{figure}
    %       \includegraphics[width=0.292\textwidth]{Coma-SDSS.jpg}
    %       \includegraphics[width=0.3\textwidth]{Coma-Rosat.pdf}
    %       \includegraphics[width=0.3\textwidth]{Coma-Planck.pdf}
    %       \vspace{-0.4cm}
    %       \caption{The Coma cluster.}
    %     \end{figure}
    \end{itemize}
\end{frame}

\begin{frame}
  \frametitle{Background: Cluster of Galaxies:}
  {\bf To understand these observational results, especially how they are formed, we need simulations.}
  % \begin{textblock}{0.25}(0.05,0.4)
  %   {\scriptsize Cluster 17 \\ credit: Gustavo Yepes.}
  % \end{textblock}
\end{frame}

\begin{frame}
  \frametitle{The Three Hundred: A successor of the {\sc nIFTy} project}
  \only<1>{
    \begin{center}
      {\Large \bf The {\sc nIFTy} galaxy cluster comparison project}\footnote{\scriptsize Ref: Sembolini et al. 2016a,b; Elahi et al. 2016; Cui et al. 2016; Arthur et al. 2017}
    \end{center}
    \alert{11} different (in both algorithms and baryon models) simulation codes are used to simulate the same galaxy cluster.
  }
  \only<2>{
    \begin{center}
      {\Large \bf What did we find? I}
    \end{center}
    {\scriptsize
    \begin{itemize}
      \item The modern SPH codes produce correct entropy profiles as AMR, moving mesh.
      \item The baryon models have larger effects than the fluid simulating techniques by mixing the entropy profiles.
    \end{itemize}}
    % \begin{figure}
    %   \vspace{-0.2cm}
    %   \includegraphics<2>[width=0.36\textwidth]{nifty-entropy.png}
    %   \includegraphics<2>[width=0.45\textwidth]{nifty-entropy-fp.png}\vspace{-0.4cm}
    %   % \caption{{\scriptsize Entropy profile. Ref: Sembolini et al. 2016a,b}}
    % \end{figure}
    % \begin{textblock}{0.12}(0.01,0.3)
    %   {\scriptsize  Entropy profile. \\ Ref: Sembolini et al. 2016a,b.}
    % \end{textblock}
  }
\end{frame}

\section{The Three Hundred}
\begin{frame}
  \frametitle{The Three Hundred: Basic information}
  \begin{itemize}
    \item The most massive ($M_{vir} > 8\times 10^{14} \hMsun$) 324 clusters are selected from the MultiDark simulation(MDPL2)\footnote{\url{https://www.cosmosim.org}}.
    \item 324 zoomed-in ICs are generated by cutting a spherical region with a radius of \alert{$15 \Mpc$} from the cluster center.
    \item[]
      \begin{table}
        \fontsize{10}{10}\selectfont
        \caption{Parameters of the Three Hundred simulations}
        \begin{tabular}{lll}
          \hline
          Parameter& Value & Description\\
          \hline
          $\Omega_M$ & 0.307 & Total Matter density parameter\\
          $\Omega_B$ & 0.048 & Baryon density parameter\\
          $\Omega_\Lambda$ & 0.693 & Cosmological Constant density parameter\\
          $h$ & 0.678  & Hubble constant in units of 100 km/s/Mpc\\
          $\sigma_8$ & 0.823 & Normalization of Power spectrum\\
          $n_s$ & 0.96  & Power index\\
          % $f_B$ & 0.0696&  Cosmic Baryon Fraction\\
          $z_{init}$ & 120 & Initial redshift of the simulations \\
          $\epsilon_{phys}$ & 6.5 & Plummer equivalent softening in $\Kpc$ \\
          % Box size & 1 $h^{-1} Gpc$ (15 $\Mpc$) & The MDPL2 simulation box on one side (the radius for each re-simulation region) \\
          Particle mass & $2.36 \ (12.7) $ & gas (dark matter) particle mass in [$10^8 \hMsun$]\\
          \hline
        \end{tabular}
      \end{table}
\end{itemize}
% \vspace{-0.6cm}
\end{frame}

\begin{frame}
  \frametitle{The Three Hundred: theoretical models}
  \begin{columns}[t]
    \begin{column}{0.55\textwidth}
      \begin{block}{hydrodynamical simulations with baryonic models:}
        {\sc Gadget-\alert{MUSIC}}: classic SPH method. Radiative cooling, star formation with both thermal and kinetic Supernove (SN) feedback.

        {\sc Gadget-\alert{X}}: modern SPH with the Wendland C4 kernel. Gas cooling with metal contributions, star formation with chemical enrichment, SN feedback with AGB phase, and AGN feedback.

        {\sc GIZMO:} running.

        {\sc Gadget-PESPH:} running.
      \end{block}
    \end{column}
    \begin{column}{0.4\textwidth}
      \begin{block}{SAMs from MultiDark-Galaxies:}
        See Alexander's talk for more details of {\sc Galacticus}, {\sc SAG} (see also sofia's talk) and {\sc SAGE} (see also Darren's talk).

        Notes: We select these catalogues from the same regions as the hydrodynamical simulations.
      \end{block}
    \end{column}
  \end{columns}
  % A coupling with MultiDark-Galaxies
\end{frame}


% \subsection{Scaling relations}
% \begin{frame}
%   \frametitle{Optical relations}
%   The complete sample is used here.
%   \begin{figure}
%     \includegraphics[width=0.33\linewidth]{Optical-relation1}
%     \includegraphics[width=0.32\linewidth]{Optical-relation2}
%     \includegraphics[width=0.32\linewidth]{Optical-relation3}
%     \vspace{-0.3cm}
%     \caption{The optical relations. Cui et al. in prep.}
%   \end{figure}
% \end{frame}
% \begin{frame}
%   \frametitle{Optical relations: the satellite stellar mass function}
%   The complete sample is used here.
%   \begin{figure}
%     \includegraphics<1>[width=0.7\linewidth]{Ssmf}
%     \vspace{-0.6cm}
%     \caption{The satellite stellar mass function. Cui et al. in prep.}
%   \end{figure}
% \end{frame}

% \begin{frame}
%   \frametitle{Gas scaling relations}
%   \begin{figure}
%     \includegraphics[width=0.48\linewidth]{T-M-relations}
%     \includegraphics[width=0.48\linewidth]{YM_relation}
%     \vspace{-0.4cm}
%     \caption{The gas relations. Cui et al. in prep.}
%   \end{figure}
% \end{frame}

\section{conclusion and future prospects}
\begin{frame}
  \frametitle{Conclusion}
  {\Large
  \begin{itemize}
    \item The baryons have a negligible impact on the halo mass for both $M_{200}$ and $M_{500}$.
    \item $\sim 20 \%$ of the complete sample is relaxed clusters, 26\% (9\%) of the sample is CC for GadgetX (MUSIC).
    \item The baryon fractions are in agreement with the observations at massive halo masses for GadgetX.
    \item The optical relations are in agreement with the observations but the galaxies from more models seem to be a little blue.
    \item The gas relations are in agreement with observations.
  \end{itemize}}
\end{frame}

\begin{frame}
  \frametitle{Future works}
  \begin{itemize}
    \item The density profiles, Mostoghiu et al. in prep.
    \item The environment effects, Wang et al. in prep.
    % \item Comparisons between different methods for estimating cluster dynamical states.
    \item The mock images in optical, X-ray and SZ (lower images). \alert{And more...}
  \end{itemize}
\end{frame}
\end{document}
